%-------------------------------------------------------------------------------
%       ABSTRACT PAGE
%-------------------------------------------------------------------------------
\addchaptertocentry{Хураангуй}
\begin{center}
{\scshape\Large \univname\par}
{\scshape\large \facname\par}\vspace{0.5cm}
{\huge\textbf{{Хураангуй}} \par}
\bigskip
{\Large{\ttitle} \par}
\bigskip
{\normalsize \shortname \par}
\addressname
\end{center}
\textit{\textbf{Түлхүүр үгс: Хиймэл оюун, VLM, LoRA, мэдлэг дамжуулалт, Зөрчлийн илрүүлэлт, Qwen загвар, хяналтын камер}}
\bigskip

Орчин үед хотын аюулгүй байдлыг хангахад видео тандалтын систем чухал үүрэг гүйцэтгэж байна. Гэвч ихэнх систем нь хүний оператороор хянагддаг тул бодит цагийн шинжилгээ хийх чадвар хязгаарлагдмал байдаг. Нэг оператор нэгэн зэрэг хянах боломжтой дэлгэцийн тоо хязгаарлагдмал бөгөөд урт хугацааны хяналтын явцад хүний анхаарал сулардаг. Энэхүү дипломын ажил нь хяналтын камерын бичлэгээс зодооны зөрчлийг автоматаар илрүүлж, Монгол хэлээр нарийвчилсан тайлбар гаргах систем боловсруулах зорилготой юм.

Судалгааны үндсэн арга нь багш--сурагч загварын мэдлэг дамжуулалтын урсгал дээр суурилна. Qwen3-VL-30B-A3B загварыг багш загвар болгон ашиглаж, Qwen2.5-VL-7B загварыг LoRA нарийн тохируулгаар сурган хөнгөн хувилбар гарган авах боломжийг судлав.

Val өгөгдөл дээрх хоёртын ангиллын загварын үнэлгээнд stillmeta/fight-video-qlora1 загвар Accuracy=0.9179, F1=0.9238, Precision=0.9238, Recall=0.9238 үр дүн үзүүлсэн. Энэ нь Qwen2-VL-7B болон Qwen2.5-VL-7B BASE загваруудаас мэдэгдэхүйц өндөр гүйцэтгэлтэй байна.

Энэхүү ажлын гол шинэлэг тал нь VLM-ийг мэдлэг дамжуулалт аргаар жижиг загварыг сургахад ашигласан, 4 түвшний зөрчлийн ангиллын системийг боловсруулсан, бодит орчинд хэрэглэх боломжтой хөнгөн загвар гарган авсан явдал юм.
